\chapter{Matching in Regular Bipartite Graphs}

\begin{definition}[Regular Graph]
	A graph \( G = (V, E) \) is called \( d \)-regular if $\forall v \in V, \deg(v) = d$.
\end{definition}

\begin{theorem}
	\label{thm:regular-matching}
	Let \( G = (L \cup R, E) \) be a \( d \)-regular bipartite graph with \( |L| = |R| = n \). Then
	\begin{itemize}
		\item \( G \) has a perfect matching.
		\item $E$ can be partitioned into $d$ perfect matchings.
	\end{itemize}
\end{theorem}
\begin{proof}[Proof 1: Using Hall's Theorem]
	If $G$ doesn't have a perfect matching, then by Hall's Theorem, there exists $S \subseteq L$ such that $|N(S)| < |S|$. However, since $G$ is $d$-regular, we have $d|S| = \sum_{v \in S} \deg(v) \le \sum_{u \in N(S)} \deg(u) = d|N(S)|$, which implies $|S| \le |N(S)|$, a contradiction. Therefore, $G$ has a perfect matching.
\end{proof}
\begin{proof}[Proof 2: Max Flow]
	Use Ford-Fulkerson algorithm to find a max flow in the graph where all capacities are 1.
	If all edges have integer capacities, the max flow we find has to be integral.

	The max flow is $n$ if and only if we can find a perfect matching.
	On the other hand, we can construct a fractional flow with $\forall e, f(e) = 1/d$ (one can easily verify it's a valid flow). Therefore, the max flow value is at least $n$, which means we can find a perfect matching.
\end{proof}

Here we present two algorithms: a deterministic one with $O(m)$ time when $d = 2^k$, and a sub-linear randomized one with $O(n \log n)$ time for general $d$.

\section{Deterministic: Gabow-Kariv Algorithm}

Here let $d = 2^k$.
We know the following theorem back in EECS 203:

\begin{theorem}
	If $\forall v\in V$, $\deg(v)$ is even, then $G$ has an Euler tour.
\end{theorem}

Unless $d = 1$, we can find an Euler tour in $G$. This gives an ordering of the edges $(e_1, e_2, e_3, \ldots, e_m)$.

For any vertex $v$, if $e_i$ goes into $v$, then $e_{i+1}$ must go out of $v$ (indices are taken modulo $m$). Therefore, if we partition the edges into odd index edges $\{e_1, e_3, e_5, \ldots\}$ and even index edges $\{e_2, e_4, e_6, \ldots\}$, then each vertex with odd in-edges must have a same number of even out-edges, and each vertex with even in-edges must have a same number of odd out-edges.

This means both the odd index edges and the even index edges form a $d/2$-regular bipartite graph. We can choose one of the edge set, and recursively apply this procedure until we get $d=1$, which gives us a perfect matching.

\begin{algorithm}[H]
	\caption{Gabow-Kariv Algorithm}
	\If{$d = 1$}{
		\Return $G$\;
	}
	Find an Euler tour $(e_1, e_2, \ldots, e_m)$ in $G$\;
	Form the graph $G_1 = (V,\{e_1, e_3, e_5, \ldots, e_{m-1}\})$\;
	\Return Gabow-Kariv($G_1$)\;
\end{algorithm}

The time complexity is $O(m + \frac{m}{2} + \frac{m}{4} + \cdots) = O(m)$.

\section{Background: Markov Chains}

\begin{definition}[Markov Chain]
	For a state space $S = \{ v_1, v_2, \ldots\}$ (finite or countable), a Markov chain is discrete-time stochastic process $X = (X_0, X_1, X_2, \ldots, X_T) in S^T$ (we can pick $X_0$ randomly or deterministically) such that $\forall t \ge 0$, $\forall i, j$, we have
	\[
		\Pr[X_{i+1} = s | X_0, X_1, \ldots, X_i = s'] = \Pr[X_{i+1} = s | X_i = s']
	\]
	(the process is \textbf{memoryless}: the history $X_0, X_1, \ldots, X_{t-1}$ is irrelevant) and
	\[
		\Pr[X_{i+1} = s | x_i = s'] = \Pr[X_{j+1} = s | X_j = s']
	\]
	(the process is \textbf{time-invariant}: the transition probability only depends on $s$ and $s'$, not on the time).

	Therefore, we can encode the transition probabilities into a matrix $P$ where

	\[
		P[a,b] = \Pr[X_{i+1} = v_b | X_i = v_a]
	\]

\end{definition}

\subsection{Random Walk on Graphs}

Define $q_i$ to be the distribution of $X_i$ if we start from $x_0=v_1$. Then
\begin{align*}
	q_0 & = (1,0,0,\ldots)  \\
	q_1 & = q_0 P           \\
	q_2 & = q_1 P = q_0 P^2 \\
	    & \vdots            \\
	q_T & = q_0 P^T
\end{align*}

For some Markov chains, $q_T$ converges to a distribution $\pi$ as $T \to \infty$. We call $\pi$ the \textbf{stationary distribution} of the Markov chain.

\begin{definition}[Stationary Distribution]
	A distribution $\pi$ is called a stationary distribution of a Markov chain with transition matrix $P$ if
	\begin{itemize}
		\item $\sum_i \pi(i)= 1$
		\item $\pi P = \pi$
	\end{itemize}
\end{definition}

For finite, irreducible Markov chains, $\pi$ is unique. However, $\lim_{n \to \infty} q_0 P^n$ may not exist when the Markov chain is periodic. For example, for a 2-state Markov chain with $P = \begin{bmatrix} 0 & 1 \\ 1 & 0 \end{bmatrix}$, we have $\pi = (1/2, 1/2)$, but $q_0 P^n$ oscillates between $(1,0)$ and $(0,1)$.

In these cases, we can instead find $\lim_{n \to \infty} q_0 (\frac{1}{2}P+\frac{1}{2}I)^n$ or $\lim_{n \to \infty} \frac{1}{n} \sum_{i=1}^{n}q_0p^i$, which both converge to $\pi$.

\subsection{Return Times}

\begin{definition}[Return Time]
	For a Markov chain $X$, $\rho(s,s'):=\E[\min\{t | X_t = s'\} | X_0 = s]$
\end{definition}

\begin{lemma}
	\[ \forall s, \pi(s) = \frac{1}{\rho(s,s)} \]
\end{lemma}
\begin{proof}[Proof Sketch]
	Let $N_n(s)$ be the number of times we visit $s$ in the first $n$ steps. Then $\lim_{n \to \infty} \frac{N_n(s)}{n} = \pi(s)$ by the definition of stationary distribution.

	Suppose the visit times are $t_0, t_1, t_2, \ldots$, and denote $Y_i = t_i - t_{i-1}$.
	Then $\E[Y_i] = \rho(s,s)$, so
	\[\lim_{k\to\infty}\frac{Y_1+Y_2+\cdots+Y_k}{k} = \rho(s,s)\]

	As $k = N_n(s)$ and $Y_1 + Y_2 + \cdots + Y_k \approx n$, we have $\lim_{n \to \infty} \frac{N_n(s)}{n} = \pi(s)$ and $\lim_{k\to\infty}\frac{n}{N_n(s)} = \rho(s,s)$, which implies $\pi(s) = \frac{1}{\rho(s,s)}$.
\end{proof}

\subsection{Eulerian Markov Chains}

\begin{definition}[Eulerian Graph]
	A directed graph $G$ is called Eulerian if $\forall v \in V, \indeg(v) = \outdeg(v)$.
\end{definition}

Consider a Markov chain defined by a Eulerian graph $G=(V,E)$ with uniform transition probabilities:

\[P[i,j] = \begin{cases}
		\frac{1}{\outdeg(v_i)} & \text{if } (v_i, v_j) \in E \\
		0                      & \text{otherwise}
	\end{cases}\]

\begin{lemma}
	The stationary distribution of the Markov chain defined by a Eulerian graph $G$ is
	\[\forall j, \pi(j) = \frac{\outdeg(v_j)}{m}.\]
\end{lemma}
\begin{proof}
	\begin{align*}
		(\pi P)[j] & = \sum_i \pi(i) P[i,j]                                                           \\
		           & = \sum_{i: (v_i, v_j) \in E} \frac{\outdeg(v_i)}{m} \cdot \frac{1}{\outdeg(v_i)} \\
		           & = \frac{\indeg(v_j)}{m}                                                          \\
		           & = \frac{\outdeg(v_j)}{m}                                                         \\
		           & = \pi(j)
	\end{align*}
\end{proof}

\section{Perfect Matching via Random Walk}

Similar to previous algorithms about matching, we try to find augmenting paths in the graph.

For a partially matched graph, we can construct a directed graph $G_M = (V', E')$ where $V' = L \cup R \cup \{s,t\}$, and $E'$ contains the following edges:
\begin{itemize}
	\item For each unmatched vertex $u \in L$, we add a directed edge $(s,u)$ to $E'$.
	\item For each unmatched vertex $v \in R$, we add a directed edge $(v,t)$ to $E'$.
	\item For each unmatched edge $(u,v)$ where $u \in L$ and $v \in R$, we add a directed edge $(u,v)$ to $E'$.
	\item For each matched edge $(u,v)$ where $u \in L$ and $v \in R$, we add a directed edge $(v,u)$ to $E'$.
\end{itemize}

\begin{lemma}
	If $W$ is a walk from $s$ to $t$ in $G_M$, then $W$ contains an augmenting path $P$ in $G$.
\end{lemma}
\begin{remark}
	$W$ itself may not be an augmenting path since it may contain cycles. If we directly flip the edges along $W$, we may end up with a matching where a vertex is matched to multiple vertices.

	However, we can trim the cycles to get an augmenting path.
\end{remark}

To analyze the probability that a random walk succeeds, we try to change the graph into a Eulerian graph.

Denote the number of free vertices in $L$ as $f$.
To start with, note that $\outdeg(s) = \indeg(t) = f$, and $\indeg(s) = \outdeg(t) = 0$. Therefore, we can contract $s$ and $t$ into a single vertex $s$, so that $\indeg(s) = \outdeg(s) = f$.

Then, for any matched vertex $u\in L$, we have $\outdeg(u) = d-1$ and $\indeg(u) = 1$. Similarly, for any matched vertex $v \in R$, we have $\outdeg(v) = 1$ and $\indeg(v) = d-1$. Therefore, we can contract all matched edges. As a result, the new vertex has $\outdeg = \indeg = d-1$.

\begin{center}
	\begin{tikzpicture}[
			free/.style={circle, draw, fill=ForestGreen!15, minimum size=0.7cm, inner sep=0pt, font=\small},
			matched/.style={circle, draw, fill=NavyBlue!15, minimum size=0.7cm, inner sep=0pt, font=\small},
			src/.style={circle, draw, fill=Gray!20, minimum size=0.7cm, inner sep=0pt, font=\small},
			>=Stealth,
		]
		% Matched pair rectangles (behind nodes)
		\draw[NavyBlue, thick, rounded corners=12pt, fill=NavyBlue!8]
		(-0.45, 3.5) rectangle (3.95, 4.5);
		\draw[NavyBlue, thick, rounded corners=12pt, fill=NavyBlue!8]
		(-0.45, 1.5) rectangle (3.95, 2.5);
		% Nodes (t relabeled s — same position, same graph)
		\node[src]     (s)  at (-2.5, -2) {$s$};
		\node[matched] (u1) at (0,     4) {$u_1$};
		\node[matched] (u2) at (0,     2) {$u_2$};
		\node[free]    (u3) at (0,     0) {$u_3$};
		\node[free]    (u4) at (0,    -2) {$u_4$};
		\node[free]    (u5) at (0,    -4) {$u_5$};
		\node[matched] (v1) at (3.5,   4) {$v_1$};
		\node[matched] (v2) at (3.5,   2) {$v_2$};
		\node[free]    (v3) at (3.5,   0) {$v_3$};
		\node[free]    (v4) at (3.5,  -2) {$v_4$};
		\node[free]    (v5) at (3.5,  -4) {$v_5$};
		\node[src]     (t)  at (6,    -2) {$s$};
		% s -> free L vertices
		\draw[->, ForestGreen, thick] (s) to[bend left=20]  (u3);
		\draw[->, ForestGreen, thick] (s) --                (u4);
		\draw[->, ForestGreen, thick] (s) to[bend right=15] (u5);
		% free R vertices -> s (formerly t)
		\draw[->, ForestGreen, thick] (v3) to[bend left=15]  (t);
		\draw[->, ForestGreen, thick] (v4) --                (t);
		\draw[->, ForestGreen, thick] (v5) to[bend right=20] (t);
		% Matched edges reversed: R -> L (inside rectangles)
		\draw[->, NavyBlue, thick] (v1) -- (u1);
		\draw[->, NavyBlue, thick] (v2) -- (u2);
		% Unmatched edges: L -> R
		\draw[->] (u1) to[bend right=8]  (v3);
		\draw[->] (u2) --                (v4);
		\draw[->] (u3) to[bend left=8]   (v1);
		\draw[->] (u3) to[bend right=8]  (v4);
		\draw[->] (u4) --                (v2);
		\draw[->] (u4) to[bend right=8]  (v5);
		\draw[->] (u5) to[bend left=5]   (v3);
		\draw[->] (u5) --                (v5);
		% Column labels
		\node[font=\small, gray] at (0,   5) {$L$};
		\node[font=\small, gray] at (3.5, 5) {$R$};
	\end{tikzpicture}
\end{center}

Finally, for any free vertex $u \in L$, we have $\outdeg(u) = d$ and $\indeg(u) = 1$, and for any free vertex $v \in R$, we have $\outdeg(v) = 1$ and $\indeg(v) = d$. We can duplicate all edges $(s,u)$ and $(v,s)$ for extra $d-1$ times, so that $\outdeg(s) = \indeg(s) = df$, $\outdeg(u) = \indeg(u) = d$, and $\outdeg(v) = \indeg(v) = d$.

\begin{center}
	\begin{tikzpicture}[
			free/.style={circle, draw, fill=ForestGreen!15, minimum size=0.7cm, inner sep=0pt, font=\small},
			contr/.style={circle, draw, fill=Violet!15, minimum size=0.7cm, inner sep=0pt, font=\small},
			src/.style={circle, draw, fill=Gray!20, minimum size=0.7cm, inner sep=0pt, font=\small},
			>=Stealth,
		]
		\node[src]   (sl) at (-2.5, -2) {$s$};
		\node[contr] (w1) at (1.75,  4) {$w_1$};
		\node[contr] (w2) at (1.75,  2) {$w_2$};
		\node[free]  (u3) at (0,     0) {$u_3$};
		\node[free]  (u4) at (0,    -2) {$u_4$};
		\node[free]  (u5) at (0,    -4) {$u_5$};
		\node[free]  (v3) at (3.5,   0) {$v_3$};
		\node[free]  (v4) at (3.5,  -2) {$v_4$};
		\node[free]  (v5) at (3.5,  -4) {$v_5$};
		\node[src]   (sr) at (6,    -2) {$s$};
		% sl -> free u: 2 parallel arrows each
		\draw[->, ForestGreen, thick] (sl) to[bend left=22]  (u3);
		\draw[->, ForestGreen, thick] (sl) to[bend left=5]   (u3);
		\draw[->, ForestGreen, thick] (sl) to[bend left=10]  (u4);
		\draw[->, ForestGreen, thick] (sl) to[bend right=10] (u4);
		\draw[->, ForestGreen, thick] (sl) to[bend right=5]  (u5);
		\draw[->, ForestGreen, thick] (sl) to[bend right=22] (u5);
		% free v -> sr: 2 parallel arrows each
		\draw[->, ForestGreen, thick] (v3) to[bend left=5]   (sr);
		\draw[->, ForestGreen, thick] (v3) to[bend left=22]  (sr);
		\draw[->, ForestGreen, thick] (v4) to[bend left=10]  (sr);
		\draw[->, ForestGreen, thick] (v4) to[bend right=10] (sr);
		\draw[->, ForestGreen, thick] (v5) to[bend right=5]  (sr);
		\draw[->, ForestGreen, thick] (v5) to[bend right=22] (sr);
		% contracted vertex edges (single copies)
		\draw[->] (u3) -- (w1);
		\draw[->] (u4) -- (w2);
		\draw[->] (w1) -- (v3);
		\draw[->] (w2) -- (v4);
		% remaining unmatched edges (single copies)
		\draw[->] (u3) to[bend right=8]  (v4);
		\draw[->] (u4) to[bend right=8]  (v5);
		\draw[->] (u5) to[bend left=5]   (v3);
		\draw[->] (u5) --                (v5);
	\end{tikzpicture}
\end{center}

This gives the full algorithm:

\begin{algorithm}[H]
	\caption{Random Walk Algorithm}
	$M \gets \emptyset$\;
	\While{$|M| < n$}{
		$G_M \gets$ the directed graph defined by $M$ where unmatched edges go from $L$ to $R$ and matched edges go from $R$ to $L$\;
		\ForEach{unmatched $u \in L$}{
			Add $d$ parallel edges from $s$ to $u$\;
		}
		\ForEach{unmatched $v \in R$}{
			Add $d$ parallel edges from $v$ to $s$\;
		}
		\ForEach{$(u,v)\in M$}{
			Contract the edge $(u,v)$\;
		}
		Do a random walk from $s$ in the new graph until we reach $s$ again\;
		Contract the walk into a simple path $P$\;
		$M \gets M \oplus P$\;
	}
	\Return $M$\;
\end{algorithm}
\begin{theorem}
	The expected running time of the random walk algorithm is $O(n \log n)$.
\end{theorem}
\begin{proof}
	$G_M$ is a Eulerian graph.
	\[\text{number of edges} = (d-1)(n-f) + 3df < nd + 2df\]
	Therefore,
	\[\rho(s,s) = \frac{\text{number of edges}}{\outdeg(s)} < \frac{nd + 2df}{df} = n/f + 2\]
	Every iteration increases the size of the matching by 1, so the expected running time is
	\[
		\sum_{f=1}^{n} (n/f + 2) = 2n + n\cdot \sum_{f=1}^{n} 1/f < 2n + n \cdot (\ln n + 1) = O(n \log n).
	\]
\end{proof}

























